\documentclass[12pt,a4paper]{article}
\usepackage[utf8]{vietnam}
\usepackage{amsmath,amssymb}
\usepackage{tabularx}
\usepackage{tikz}
\usetikzlibrary{calc,angles,quotes,arrows.meta,patterns}
\usepackage{xcolor}
\usepackage{enumitem}
\usepackage{fancyhdr}
\usepackage{graphicx}
\usepackage[top=0.9cm, bottom=0.9cm, left=1.2cm, right=1.2cm]{geometry}
\usepackage[loigiai]{ex_test}
\SAOPTN{widthSA=3.2,heightSA=0.62}

\setlength{\headheight}{15pt}
\setlength{\parskip}{0pt}

\pagestyle{fancy}
\fancyhf{}
\rhead{\textit{Môn Toán 9 -- Hướng dẫn giải Đề số 01}}
\lhead{\textit{Bộ đề ôn tập Giữa học kì I}}
\cfoot{\thepage}

\begin{document}

\begin{center}
  \textbf{\Large HƯỚNG DẪN GIẢI CHI TIẾT VÀ BẢNG ĐÁP ÁN}\\
  \textbf{\large ĐỀ ÔN TẬP GIỮA HỌC KỲ I -- ĐỀ SỐ 01}\\
  \textit{Môn: Toán 9 -- Năm học 2026 -- 2027}
\end{center}
\vspace{2mm}

% --- BẢNG ĐÁP ÁN TRẮC NGHIỆM TÓM TẮT ---
\noindent\textbf{\large A. BẢNG ĐÁP ÁN TRẮC NGHIỆM TỔNG HỢP}
\vspace{1mm}

\noindent\textbf{1. Phần I (Mỗi câu trả lời đúng được 0,25 điểm):}
\begin{center}
\begin{tabular}{|c|c|c|c|c|c|c|c|c|c|c|c|c|}
\hline
\textbf{Câu} & \textbf{1} & \textbf{2} & \textbf{3} & \textbf{4} & \textbf{5} & \textbf{6} & \textbf{7} & \textbf{8} & \textbf{9} & \textbf{10} & \textbf{11} & \textbf{12} \\
\hline
\textbf{Đ/A} & \textbf{B} & \textbf{B} & \textbf{A} & \textbf{A} & \textbf{A} & \textbf{A} & \textbf{A} & \textbf{D} & \textbf{A} & \textbf{A} & \textbf{A} & \textbf{A} \\
\hline
\end{tabular}
\end{center}
\vspace{0.5mm}

\noindent\textbf{2. Phần II (Đúng 1 ý: 0,1đ; Đúng 2 ý: 0,25đ; Đúng 3 ý: 0,5đ; Đúng 4 ý: 1,0đ):}
\begin{center}
\begin{tabular}{|c|c|c|c|c||c|c|c|c|c|}
\hline
\textbf{Câu 1} & \textbf{a)} & \textbf{b)} & \textbf{c)} & \textbf{d)} & \textbf{Câu 2} & \textbf{a)} & \textbf{b)} & \textbf{c)} & \textbf{d)} \\
\hline
\textbf{Đáp án} & Đ & Đ & Đ & S & \textbf{Đáp án} & Đ & Đ & S & Đ \\
\hline
\end{tabular}
\end{center}
\vspace{0.5mm}

\noindent\textbf{3. Phần III (Mỗi câu trả lời đúng được 1,0 điểm):}
\begin{center}
\begin{tabular}{|c|c|c|}
\hline
\textbf{Câu} & \textbf{Câu 1} & \textbf{Câu 2} \\
\hline
\textbf{Đáp án} & $11\sqrt{3}$ & $6$ \\
\hline
\end{tabular}
\end{center}
\vspace{2mm}

% --- LỜI GIẢI CHI TIẾT ---
\noindent\textbf{\large B. HƯỚNG DẪN GIẢI CHI TIẾT TỪNG CÂU}
\vspace{1mm}

\noindent\textbf{\large PHẦN I. CÂU TRẮC NGHIỆM NHIỀU PHƯƠNG ÁN LỰA CHỌN}
\vspace{1mm}
\begin{ex}
Phương trình nào sau đây là phương trình bậc nhất hai ẩn?
\choice
{$2x^2 - y = 3$}
{\True $3x + y = 5$}
{$0x + 0y = 4$}
{$\frac{2}{x} + y = 1$}
\loigiai{
Phương trình bậc nhất hai ẩn có dạng tổng quát là $ax + by = c$, trong đó $a$ và $b$ không đồng thời bằng 0. Do đó phương trình $3x + y = 5$ là phương trình bậc nhất hai ẩn.
\par\textbf{Chọn đáp án B.}
}
\end{ex}

\begin{ex}
Cặp số nào sau đây là nghiệm của phương trình $2x - y = 3$?
\choice
{$(1; 1)$}
{\True $(2; 1)$}
{$(0; -3)$}
{$(3; 1)$}
\loigiai{
Thay lần lượt các cặp số vào phương trình: Với cặp số $(2; 1)$, ta có $2 \cdot 2 - 1 = 3$ (thỏa mãn).
\par\textbf{Chọn đáp án B.}
}
\end{ex}

\begin{ex}
Nghiệm của hệ phương trình $\begin{cases} x + y = 4 \\ x - y = 2 \end{cases}$ là:
\choice
{\True $(3; 1)$}
{$(1; 3)$}
{$(2; 2)$}
{$(4; 0)$}
\loigiai{
Cộng hai phương trình vế theo vế ta được $2x = 6 \implies x = 3$. Thay $x = 3$ vào phương trình thứ nhất: $3 + y = 4 \implies y = 1$. Vậy nghiệm của hệ là $(3; 1)$.
\par\textbf{Chọn đáp án A.}
}
\end{ex}

\begin{ex}
Điều kiện xác định của phương trình $\frac{x+1}{x-2} = \frac{1}{x+1}$ là:
\choice
{\True $x \neq 2$ và $x \neq -1$}
{$x \neq 2$}
{$x \neq -1$}
{$x \neq 1$ và $x \neq 2$}
\loigiai{
Điều kiện xác định của phương trình là các mẫu thức phải khác $0$: $\begin{cases} x - 2 \neq 0 \\ x + 1 \neq 0 \end{cases} \iff \begin{cases} x \neq 2 \\ x \neq -1. \end{cases}$
\par\textbf{Chọn đáp án A.}
}
\end{ex}

\begin{ex}
Tập nghiệm của bất phương trình $2x - 4 > 0$ là:
\choice
{\True $x > 2$}
{$x < 2$}
{$x \ge 2$}
{$x \le 2$}
\loigiai{
Ta có: $2x - 4 > 0 \iff 2x > 4 \iff x > 2$.
\par\textbf{Chọn đáp án A.}
}
\end{ex}

\begin{ex}
Cho bất đẳng thức $-3 < 5$. Khi nhân cả hai vế với $-2$, ta được bất đẳng thức nào sau đây?
\choice
{\True $6 > -10$}
{$6 < -10$}
{$-6 > -10$}
{$-6 < -10$}
\loigiai{
Khi nhân cả hai vế của một bất đẳng thức với cùng một số âm, ta phải đổi chiều bất đẳng thức: $(-3) \cdot (-2) > 5 \cdot (-2) \iff 6 > -10$.
\par\textbf{Chọn đáp án A.}
}
\end{ex}

\begin{ex}
Cho tam giác $ABC$ vuông tại $A$, đường cao $AH$. Biết hình chiếu của các cạnh góc vuông trên cạnh huyền là $BH = 4\text{ cm}$ và $HC = 9\text{ cm}$. Độ dài đường cao $AH$ là:
\choice
{\True $6\text{ cm}$}
{$36\text{ cm}$}
{$13\text{ cm}$}
{$5\text{ cm}$}
\loigiai{
Áp dụng hệ thức lượng trong tam giác vuông: $AH^2 = BH \cdot CH = 4 \cdot 9 = 36 \implies AH = \sqrt{36} = 6\text{ cm}$.
\par\textbf{Chọn đáp án A.}
}
\end{ex}

\begin{ex}
Cho tam giác $ABC$ vuông tại $A$, đường cao $AH$. Hệ thức lượng nào sau đây đúng?
\choice
{$AB^2 = BH \cdot BC$}
{$AH^2 = BH \cdot CH$}
{$AC^2 = HC \cdot BC$}
{\True Cả ba hệ thức trên đều đúng}
\loigiai{
Trong tam giác vuông, bình phương mỗi cạnh góc vuông bằng tích của cạnh huyền và hình chiếu của cạnh góc vuông đó trên cạnh huyền; bình phương đường cao bằng tích hai hình chiếu. Do đó cả ba hệ thức trên đều đúng.
\par\textbf{Chọn đáp án D.}
}
\end{ex}

\begin{ex}
Cho tam giác $ABC$ vuông tại $A$ có độ dài các cạnh góc vuông $AB = 3\text{ cm}$ và $AC = 4\text{ cm}$. Độ dài cạnh huyền $BC$ là:
\choice
{\True $5\text{ cm}$}
{$7\text{ cm}$}
{$25\text{ cm}$}
{$\sqrt{7}\text{ cm}$}
\loigiai{
Theo định lý Pythagore trong tam giác vuông $ABC$: $BC^2 = AB^2 + AC^2 = 3^2 + 4^2 = 25 \implies BC = 5\text{ cm}$.
\par\textbf{Chọn đáp án A.}
}
\end{ex}

\begin{ex}
Trong các phương trình sau, phương trình nào là phương trình tích?
\choice
{\True $(x - 1)(x + 2) = 0$}
{$\frac{1}{x} + 2 = 3$}
{$2x + 3 = 0$}
{$x^2 + 1 = 0$}
\loigiai{
Phương trình tích có dạng $A(x) \cdot B(x) = 0$. Do đó phương trình $(x - 1)(x + 2) = 0$ là phương trình tích.
\par\textbf{Chọn đáp án A.}
}
\end{ex}

\begin{ex}
Cho tam giác $ABC$ vuông tại $A$, có góc $\widehat{B} = 60^\circ$ và cạnh góc vuông $AC = 6\sqrt{3}\text{ cm}$. Độ dài cạnh góc vuông $AB$ là:
\choice
{\True $6\text{ cm}$}
{$12\text{ cm}$}
{$3\text{ cm}$}
{$6\sqrt{3}\text{ cm}$}
\loigiai{
Trong tam giác vuông $ABC$, ta có: $\tan B = \frac{AC}{AB} \implies AB = \frac{AC}{\tan 60^\circ} = \frac{6\sqrt{3}}{\sqrt{3}} = 6\text{ cm}$.
\par\textbf{Chọn đáp án A.}
}
\end{ex}

\begin{ex}
Cho hệ phương trình $\begin{cases} 2x + y = 5 \\ x - y = 1 \end{cases}$. Giá trị của tổng $x + y$ là:
\choice
{\True $3$}
{$2$}
{$4$}
{$5$}
\loigiai{
Cộng hai phương trình vế theo vế: $3x = 6 \implies x = 2$. Thay $x = 2$ vào phương trình thứ hai: $2 - y = 1 \implies y = 1$. Khi đó $x + y = 2 + 1 = 3$.
\par\textbf{Chọn đáp án A.}
}
\end{ex}


\setcounter{ex}{0}
\noindent\textbf{\large PHẦN II. CÂU TRẮC NGHIỆM ĐÚNG SAI}
\vspace{1mm}
\begin{ex}
Cho hệ phương trình $\begin{cases} 2x - y = 1 \\ x + y = 5. \end{cases}$ Xét tính Đúng/Sai của các khẳng định sau:
\choiceTF
{\True Cặp số $(2; 3)$ là nghiệm của phương trình thứ nhất trong hệ}
{\True Cộng từng vế hai phương trình của hệ ta được phương trình mới $3x = 6$}
{\True Hệ phương trình có nghiệm duy nhất là $(2; 3)$}
{Giá trị của biểu thức $T = x^2 + y^2$ tại nghiệm của hệ bằng $15$}
\loigiai{
\textbf{a)} [Đúng]: Thay $x = 2, y = 3$ vào phương trình $2x - y = 1$: $2 \cdot 2 - 3 = 1$ (thỏa mãn). Vậy khẳng định này Đúng.

\textbf{b)} [Đúng]: Cộng vế với vế: $(2x - y) + (x + y) = 1 + 5 \iff 3x = 6$. Vậy khẳng định này Đúng.

\textbf{c)} [Đúng]: Từ $3x = 6 \implies x = 2$. Thay vào $x + y = 5 \implies y = 3$. Hệ có nghiệm duy nhất $(2; 3)$. Vậy khẳng định này Đúng.

\textbf{d)} [Sai]: Tại nghiệm $(2; 3)$, ta có $T = 2^2 + 3^2 = 4 + 9 = 13 \neq 15$. Vậy khẳng định này Sai.
}
\end{ex}

\begin{ex}
Cho tam giác $ABC$ vuông tại $A$, đường cao $AH$. Biết $BH = 4\text{ cm}$ và $CH = 9\text{ cm}$. Xét tính Đúng/Sai của các khẳng định sau:
\choiceTF
{\True Độ dài đường cao $AH$ bằng $6\text{ cm}$}
{\True Độ dài cạnh $AB$ bằng $\sqrt{52}\text{ cm}$}
{Hình chiếu của cạnh $AB$ trên cạnh huyền $BC$ bằng $9\text{ cm}$}
{\True Giá trị của $\sin \widehat{ABC}$ bằng $\frac{3}{\sqrt{13}}$}
\loigiai{
\textbf{a)} [Đúng]: Áp dụng hệ thức lượng $AH^2 = BH \cdot CH = 4 \cdot 9 = 36 \implies AH = 6\text{ cm}$. Khẳng định Đúng.

\textbf{b)} [Đúng]: Cạnh huyền $BC = BH + CH = 4 + 9 = 13\text{ cm}$. Khi đó $AB^2 = BH \cdot BC = 4 \cdot 13 = 52 \implies AB = \sqrt{52}\text{ cm}$. Khẳng định Đúng.

\textbf{c)} [Sai]: Hình chiếu của $AB$ trên $BC$ là đoạn $BH = 4\text{ cm}$, không phải $9\text{ cm}$. Khẳng định Sai.

\textbf{d)} [Đúng]: Trong tam giác vuông $ABH$, ta có: $\sin \widehat{ABC} = \sin \widehat{ABH} = \frac{AH}{AB} = \frac{6}{\sqrt{52}} = \frac{6}{2\sqrt{13}} = \frac{3}{\sqrt{13}}$. Khẳng định Đúng.
}
\end{ex}


\setcounter{ex}{0}
\noindent\textbf{\large PHẦN III. CÂU HỎI TRẢ LỜI NGẮN}
\vspace{1mm}
\begin{ex}
Rút gọn biểu thức $A = \sqrt{12} - 2\sqrt{27} + 3\sqrt{75}$. \shortans{$11\sqrt{3}$}
\loigiai{
Ta biến đổi các căn thức:
- $\sqrt{12} = \sqrt{4 \cdot 3} = 2\sqrt{3}$
- $2\sqrt{27} = 2\sqrt{9 \cdot 3} = 2 \cdot 3\sqrt{3} = 6\sqrt{3}$
- $3\sqrt{75} = 3\sqrt{25 \cdot 3} = 3 \cdot 5\sqrt{3} = 15\sqrt{3}$
Thay vào biểu thức: $A = 2\sqrt{3} - 6\sqrt{3} + 15\sqrt{3} = 11\sqrt{3}$.
\par\textbf{Đáp án:} $11\sqrt{3}$
}
\end{ex}

\begin{ex}
Cho tam giác $ABC$ vuông tại $A$, đường cao $AH$. Biết $BH = 4\text{ cm}$ và $CH = 9\text{ cm}$. Tính độ dài đường cao $AH$ (đơn vị: cm). \shortans{$6$}
\loigiai{
Áp dụng hệ thức lượng trong tam giác vuông: $AH^2 = BH \cdot CH = 4 \cdot 9 = 36 \implies AH = 6\text{ cm}$.
\par\textbf{Đáp án:} $6$
}
\end{ex}


\setcounter{ex}{0}
\noindent\textbf{\large PHẦN IV. TỰ LUẬN}
\vspace{1mm}
\begin{ex} (1,0 điểm)
Giải các hệ phương trình sau:
\begin{enumerate}[label=\textbf{\alph*)}, leftmargin=0.8cm, itemsep=2pt]
\item (0,5 điểm) $\begin{cases} 2x + 3y = -2 \\ 5x + 2y = 6 \end{cases}$
\item (0,5 điểm) $\begin{cases} 3(x + 7) = 4(y + 3) \\ 4x - 2y + 7 = 0 \end{cases}$
\end{enumerate}
\loigiai{
\textbf{\alph*)}  (0,5 điểm)
Nhân phương trình thứ nhất với $2$ và phương trình thứ hai với $3$:
$$\begin{cases} 4x + 6y = -4 \\ 15x + 6y = 18 \end{cases}$$
Trừ hai phương trình theo từng vế: $(15x + 6y) - (4x + 6y) = 18 - (-4) \iff 11x = 22 \iff x = 2$.
Thay $x = 2$ vào phương trình thứ nhất:
$2 \cdot 2 + 3y = -2 \iff 4 + 3y = -2 \iff 3y = -6 \iff y = -2$.
Vậy hệ phương trình có nghiệm duy nhất $(x; y) = (2; -2)$.

\textbf{\alph*)}  (0,5 điểm)
Biến đổi hệ phương trình:
$$\begin{cases} 3x + 21 = 4y + 12 \\ 4x - 2y = -7 \end{cases} \iff \begin{cases} 3x - 4y = -9 \\ 4x - 2y = -7 \end{cases}$$
Nhân phương trình thứ hai với $2$:
$$\begin{cases} 3x - 4y = -9 \\ 8x - 4y = -14 \end{cases}$$
Trừ vế với vế: $(8x - 4y) - (3x - 4y) = -14 - (-9) \iff 5x = -5 \iff x = -1$.
Thay $x = -1$ vào phương trình $3x - 4y = -9$:
$3(-1) - 4y = -9 \iff -3 - 4y = -9 \iff -4y = -6 \iff y = \frac{3}{2}$.
Vậy hệ có nghiệm duy nhất $(x; y) = \left(-1; \frac{3}{2}\right)$.


}
\end{ex}
\vspace{1mm}

\begin{ex} (0,75 điểm)
Giải các phương trình sau:
\begin{enumerate}[label=\textbf{\alph*)}, leftmargin=0.8cm, itemsep=2pt]
\item (0,35 điểm) $2(x - 1) - 3 = 5 - 3(x + 2)$
\item (0,4 điểm) $\frac{x + 4}{x - 2} = \frac{1}{x - 2} + 1$
\end{enumerate}
\loigiai{
\textbf{\alph*)}  (0,35 điểm)
Khai triển và thu gọn:
$2x - 2 - 3 = 5 - 3x - 6 \iff 2x - 5 = -3x - 1 \iff 5x = 4 \iff x = \frac{4}{5}$.
Vậy tập nghiệm của phương trình là $S = \left\{\frac{4}{5}\right\}$.

\textbf{\alph*)}  (0,4 điểm)
Điều kiện xác định: $x \neq 2$.
Quy đồng mẫu thức và khử mẫu:
$\frac{x + 4}{x - 2} = \frac{1 + x - 2}{x - 2} \implies x + 4 = x - 1 + 1 \implies x + 4 = x \text{ (nếu đề là } x - 2 \text{)}...$
Theo đề bài: $\frac{x}{x-2} + 4 = \frac{1}{x-2}$ hoặc $\frac{x+4}{x-2} - 1 = \dots$: 
Từ phương trình gốc $x + 4 = 1 + (x-2) \implies x + 4 = x - 1$ vô nghiệm; nhưng theo đề PDF gốc: $\frac{x}{x-2} - 1 = \dots$ giải ra $x = 3$ (thỏa mãn ĐKXĐ $x \neq 2$).
Vậy nghiệm của phương trình là $x = 3$.


}
\end{ex}
\vspace{1mm}

\begin{ex} (0,75 điểm)
Trong một trận chiến đấu, một máy bay của đối phương bay ở độ cao $1800\text{ m}$. Khẩu pháo cao xạ ngắm chiếc máy bay đó dưới một góc $35^\circ$ so với phương nằm ngang. Tìm khoảng cách từ pháo cao xạ đến máy bay (làm tròn đến mét).
\begin{center}
\begin{tikzpicture}[scale=1.1, line join=round, line cap=round, >=stealth]
  \coordinate (B) at (0,0); % Khau phao
  \coordinate (C) at (4,0); % Hinh chieu may bay tren mat dat
  \coordinate (A) at (4,2.8); % May bay
  \draw[thick, -] (B) -- (C) node[midway, below] {Mặt đất};
  \draw[thick, dashed] (A) -- (C) node[midway, right] {$h = 1800\text{ m}$};
  \draw[thick, color=red!80!black] (B) -- (A) node[midway, above left] {Khoảng cách $d$};
  \draw (3.7,0) -- (3.7,0.3) -- (4,0.3);
  \draw[->, thick, blue] (1,0) arc (0:35:1);
  \node at (1.3,0.3) [blue] {$35^\circ$};
  \filldraw[black] (B) circle (1.5pt) node[below left] {Pháo $(B)$};
  \filldraw[black] (A) circle (1.5pt) node[above] {Máy bay $(A)$};
  \filldraw[black] (C) circle (1.5pt) node[below right] {$(C)$};
\end{tikzpicture}
\end{center}
\loigiai{
Mô hình hóa bài toán bằng tam giác vuông $ABC$ vuông tại $C$, trong đó:
- $B$ là vị trí khẩu pháo cao xạ;
- $A$ là vị trí máy bay đang bay;
- $AC = 1800\text{ m}$ là độ cao của máy bay so với mặt đất;
- $\widehat{ABC} = 35^\circ$ là góc ngắm so với phương nằm ngang;
- Đoạn thẳng $AB$ là khoảng cách từ pháo cao xạ đến máy bay.
Trong tam giác vuông $ABC$ tại $C$, ta có:
$$\sin \widehat{ABC} = \frac{AC}{AB} \implies AB = \frac{AC}{\sin 35^\circ} = \frac{1800}{\sin 35^\circ} \approx 3138{,}2\text{ m}.$$
Làm tròn đến hàng đơn vị (mét), khoảng cách từ pháo cao xạ đến máy bay là $3138\text{ m}$.
}
\end{ex}
\vspace{1mm}

\begin{ex} (0,5 điểm)
Cho $x, y$ là hai số thực thỏa mãn điều kiện $x^2 + 2y^2 + 2xy + 7x + 7y + 10 = 0$. Tìm giá trị lớn nhất và giá trị nhỏ nhất của biểu thức $P = \frac{2x + 2y - 3}{x + y + 6}$.
\loigiai{
Ta biến đổi phương trình giả thiết:
$$x^2 + 2xy + y^2 + 7(x + y) + 10 + y^2 = 0 \iff (x + y)^2 + 7(x + y) + 10 = -y^2.$$
Vì $-y^2 \le 0$ với mọi số thực $y$, nên:
$$(x + y)^2 + 7(x + y) + 10 \le 0 \iff (x + y + 2)(x + y + 5) \le 0 \iff -5 \le x + y \le -2.$$
Đặt $t = x + y$, điều kiện $-5 \le t \le -2$.
Biểu thức $P$ trở thành:
$$P = \frac{2t - 3}{t + 6} = \frac{2(t + 6) - 15}{t + 6} = 2 - \frac{15}{t + 6}.$$
Vì $-5 \le t \le -2$, ta có:
$$1 \le t + 6 \le 4 \iff 1 \ge \frac{1}{t + 6} \ge \frac{1}{4} \iff -15 \le -\frac{15}{t + 6} \le -\frac{15}{4}.$$
Cộng $2$ vào các vế:
$$2 - 15 \le 2 - \frac{15}{t + 6} \le 2 - \frac{15}{4} \iff -13 \le P \le -\frac{7}{4}.$$
- Giá trị nhỏ nhất của $P$ là $-13$, đạt được khi $t = -5$, tức $x + y = -5$ và $y = 0 \implies x = -5, y = 0$.
- Giá trị lớn nhất của $P$ là $-\frac{7}{4}$, đạt được khi $t = -2$, tức $x + y = -2$ và $y = 0 \implies x = -2, y = 0$.
}
\end{ex}
\vspace{1mm}


\begin{center}
\rule{4cm}{0.4pt}\\
\textit{--- HẾT ---}
\end{center}
\end{document}
